% Editable source for MadFieldRunner_User_Guide.pdf.
% Rebuild here with latexmk -pdf -interaction=nonstopmode -halt-on-error MadFieldRunner_User_Guide.tex
\documentclass[11pt,letterpaper]{article}
\usepackage[T1]{fontenc}
\usepackage{lmodern}
\usepackage[margin=0.78in,top=0.68in,bottom=0.67in,footskip=24pt]{geometry}
\usepackage{microtype}
\usepackage{xcolor}
\usepackage{array,booktabs,tabularx}
\usepackage{titlesec}
\usepackage{listings}
\usepackage{fancyhdr}
\usepackage[hidelinks]{hyperref}

\definecolor{ink}{HTML}{263846}
\definecolor{navy}{HTML}{193348}
\definecolor{muted}{HTML}{526879}
\definecolor{teal}{HTML}{147078}
\definecolor{pale}{HTML}{F0F5F7}
\definecolor{line}{HTML}{D5E1E7}

\renewcommand{\familydefault}{\sfdefault}
\color{ink}
\setlength{\parindent}{0pt}
\setlength{\parskip}{4pt}
\renewcommand{\arraystretch}{1.21}
\DisableLigatures{encoding = T1, family = lmtt}
\newcommand{\code}[1]{\texttt{\detokenize{#1}}}
\newcolumntype{Y}{>{\raggedright\arraybackslash}X}
\titleformat{\section}[hang]{\large\bfseries\color{navy}}{\color{teal}\thesection.}{0.45em}{}
\titlespacing*{\section}{0pt}{12pt plus 2pt}{4pt}
\hypersetup{pdftitle={MadFieldRunner User Guide},pdfsubject={Terminal workflow and saved results}}
\pagestyle{fancy}
\fancyhf{}
\fancyfoot[L]{\footnotesize\color{muted}MadFieldRunner User Guide}
\fancyfoot[R]{\footnotesize\color{muted}\thepage}
\renewcommand{\headrulewidth}{0pt}
\lstset{
  basicstyle=\ttfamily\small\color{navy},
  backgroundcolor=\color{pale},rulecolor=\color{line},
  frame=single,framerule=0.4pt,framesep=5pt,xleftmargin=5pt,xrightmargin=5pt,
  aboveskip=7pt,belowskip=6pt,columns=fullflexible,keepspaces=true,
  showstringspaces=false,upquote=true
}

\begin{document}

{\Huge\bfseries\color{navy}MadFieldRunner\par}
\vspace{2pt}
{\large\color{muted}Run the sextet example and check your result\par}

\section{Prepare your terminal}
Download and extract the \href{https://github.com/HengamehB/MadFieldRunner/archive/refs/heads/main.zip}{\textcolor{teal}{source ZIP}}
or clone the \href{https://github.com/HengamehB/MadFieldRunner}{\textcolor{teal}{repository}}.
Run the commands from the package folder. Use Bash on macOS or Linux. On
Windows, use an Ubuntu terminal in WSL2 with a Linux Python and a Linux
Wolfram kernel. A working, licensed Mathematica or Wolfram Engine in that
environment satisfies the Wolfram requirement. Install and activate a kernel
only if one is missing. The package includes FeynRules and MadGraph.

Activate your Python 3.11 environment with NumPy and SciPy and run the
check below. If you need a new Conda environment, use
\code{conda create -n madfieldrunner -c conda-forge python=3.11 numpy scipy}
and then \code{conda activate madfieldrunner}.

\begin{lstlisting}
./run.sh check
\end{lstlisting}
If the check passes, continue to step~2. Install any missing prerequisites
it identifies, including GNU Fortran, a C++ compiler, Make or Perl. If a
tool is installed but its path was missed, run \code{./run.sh configure-tools}.
Enter \code{edit} to correct the path or \code{y} to save the detected paths
in \code{.config}, then repeat \code{./run.sh check}.

\section{Run the supplied model}
The two files in \code{input/} have distinct roles. \code{sample_particle_sextet.txt}
defines the scalar colour sextet and triplet. The coefficients and calculation
settings are in \code{sample_model_sextet.txt}. Run the model directly.
\begin{lstlisting}
./run.sh rates input/sample_model_sextet.txt --selected C0004 --page-size 5
\end{lstlisting}
This exports the UFO, calculates two-body decays and integrates the
$pp\to x_1$ production channel. The example uses $13.6$~TeV, the
\code{nn23lo1} PDF, dynamical scale choice 3 and seed 1.
The terminal shows five rows. The saved CSV contains all 1,925 histories.

\section{Compare with the supplied result}
Open \code{report.html} in the folder printed after \textbf{Run files}.
Compare it with \code{sample_output/report.html}, or check the saved
tables with this command. Replace the quoted path with your run folder.
\begin{lstlisting}
python sample_output/compare.py "path/to/saved run"
\end{lstlisting}
For a manual check, find \code{C0004} and $x_1\to u\,u$ in \code{ranking.csv}.

\begin{tabularx}{\linewidth}{@{}Yl@{}}
\toprule
\textbf{Quantity} & \textbf{Reference value} \\
\midrule
Production cross section for $pp\to x_1$ & $1.1691955 \pm 0.0045663$ pb \\
Branching fraction for $x_1\to u\,u$ & $0.3392486$ \\
Ranked rate for $pp\to x_1\to u\,u$ & $0.39664794 \pm 0.00154912$ pb \\
\bottomrule
\end{tabularx}

Thirty histories have numerical rates. The other 1,895 have blank rates
because their production channels were not integrated. The errors above
come from Monte Carlo integration. For a numerical cross-check, compare
the two production cross sections within three times their errors added
in quadrature.

\newpage
{\LARGE\bfseries\color{navy}Use your own filenames and settings\par}

\section{Choose a model filename}
\textbf{To keep the sample coefficients}, copy the supplied model to a new
name, then edit it. Filenames may contain spaces when quoted.

\begin{lstlisting}
cp input/sample_model_sextet.txt "input/my scan.txt"
./run.sh validate "input/my scan.txt"
\end{lstlisting}

\textbf{To start from a different particle definition}, copy the particle
file under a new name. Edit \code{X} and \code{Y} in your copy first
(use \code{Y = None} for a single new multiplet). Then run
\code{init} and use \code{--model} to choose a fresh destination.

\begin{lstlisting}
cp input/sample_particle_sextet.txt "input/my particles.txt"
# Edit "input/my particles.txt" before running init.
./run.sh init "input/my particles.txt" --model "input/new model.txt"
./run.sh inputs "input/new model.txt"
\end{lstlisting}

For \code{sample_particle_sextet.txt}, the paired default is
\code{sample_model_sextet.txt}. \code{init} enables diagonal kinetic and mass terms. It \textbf{replaces
its destination without a backup}, so choose a different path from any tuned
model you want to keep. In the new model, set operator switches and
coefficients. \code{inputs} explains units and array ordering. Declare
physically stable states in \code{stable_particles}. Validation derives
the spectrum with an unmixed 125~GeV Higgs and zero new-scalar VEVs.

Keep the model's \code{_metadata} line and the adjacent
\code{.model_metadata/} folder together when moving it. Pass the model path
to each command. To make it the default, set \code{model = input/my scan.txt}
under \code{[project]} in \code{.config}.

\section{Choose what to calculate}
\begin{lstlisting}
./run.sh channels "input/my scan.txt"
./run.sh histories "input/my scan.txt" --page-size 20
./run.sh run "input/my scan.txt"
\end{lstlisting}
\code{channels} lists the production IDs for your model. Use \code{rates}
with \code{--selected} and one of those IDs for a selected integration.
\code{histories} calculates widths and builds cascades before any production
integration. \code{--page-size} limits the display only. \code{run} with
no selection integrates every discovered channel. Press \textbf{Ctrl+C} to
stop and repeat your command to resume verified work.

\section{Locate saved tables and refresh the report}
The example writes to \code{output/<model_name>/<run-id>/}. Set
\code{model_name} to label your own run. A relative \code{output} in the
model starts beside that file. In the run folder, \code{cross_sections.csv}
contains the production integrals and errors. \code{total_widths.csv} and
\code{decays.csv} contain widths and branching fractions. In
\code{manifest.json}, \code{requested_complete} records completion of the
requested selection and \code{all_rates_complete} covers all channels.

\begin{lstlisting}
./run.sh report --run-dir "path/to/saved run"
\end{lstlisting}
This rebuilds the HTML from saved tables. The ranking uses tree-level
production and narrow-width cascades. Loop-induced production is excluded,
omitted decay contributions can be uncertain, and overlapping rows must
not be summed. Scale choice 3 depends on external particles, so its
numerical scale may differ in an explicit decayed process.

\end{document}
